%% kapa: exam=1 round=30 year=2019 period=1 form=A subject=economics subject_ko=경제학원론 seq=15 number=55
%% source: Q-Net 감정평가사 기출문제 (cst00309) artlSeq=5205205 file=extracted/1008/5205205/2019년 제30회 감정평가사 1차시험 문제지/2019년 제30회 감정평가사 1차시험 문제지/2019년도 제30회 감정평가사 1차 1교시 A형.hwp
%% conversion: kordoc 4.24.51 -> tools/gen1.py (auto); review=proofread-sonnet; flags=none
\begin{question}{55}
두 공장 1, 2를 운영하고 있는 기업 A의 비용함수는 각각 ${ C_{ 1 } (q_{ 1 } )=q_{ 1 }^{ 2 } , C_{ 2 } (q_{ 2 } )=2q_{ 2 } }$이다. 총비용을 최소화하여 5단위를 생산하는 경우, 공장 1, 2에서의 생산량은? (단, ${ q_{ 1 } }$은 공장 1의 생산량, ${ q_{ 2 } }$는 공장 2의 생산량이다.)
\begin{choices}
  \item $q_{ 1 } =5, q_{ 2 } =0$
  \item $q_{ 1 } =4, q_{ 2 } =1$
  \item $q_{ 1 } =3, q_{ 2 } =2$
  \item $q_{ 1 } =2, q_{ 2 } =3$
  \item $q_{ 1 } =1, q_{ 2 } =4$
\end{choices}
\end{question}
\answer{5}
